

\begin{table}
\centering
\begin{tabular}{| c | c | c |}
\hline
\# of observed events	                        	& Rate only      & Rate + Shape		    	\\
\hline
$N_{obs} = 2$      		                    	& 0.510          &          	\\
$ N_{obs} = 3$      		                    	& 0.251          &          	\\
$ \bf N_{obs} = 4$      		           	         & \bf 0.101     &             \\
$N_{obs} = 5$      		                    	& 0.0343        &        	\\
$N_{obs} = 6$      		                    	& 0.0103        &        	\\
\hline
\end{tabular}
\caption{
Summary of the expected rate-only p-values for different number of 
observed events. The simulated number of POT corresponds to Run 1-6.
\NumToysLarge toys are used to get the distribution of the test statistic and the 
Asimov data is used to compute the p-value.
The Asimov dataset for $\beta=1$ is in bold.
}
\label{tab:summary_sensi_rateonly}
\end{table}




\begin{table}
\centering
\begin{tabular}{| c | c | c |}
\hline
P-value	                                          	& Asimov    	& Mean 		    	\\
\hline
Rate-only                          			& 0.109           	& 0.269                  	\\
Rate+shape                          		& 0.0858          	& 0.170                 	\\
\hline
\end{tabular}
\caption{
Summary of the expected rate-only and rate+shape p-values for 
the Run 1-6 statistic.
The Asimov p-value is calculated using the Asimov dataset while the 
mean p-value is calculated making \NumToysLarge toys and taking the mean 
of the p-value distribution.
}
\label{tab:summary_sensi_lowpot}
\end{table}




%% Comparison with many toys
%
%\begin{table}
%\centering
%\begin{tabular}{| c | c | c |}
%\hline
%P-value	                                          			& Asimov    	  & Mean 		    	\\
%\hline
%Rate-only (\NumToys toys)         			& 0.110           & 0.207                  	\\
%Rate+shape (\NumToys toys)              		& 0.0726          & 0.143                  	\\
%\hline
%Rate-only (\NumToysLarge toys)        		& 0.109           & 0.203                	\\
%Rate+shape (\NumToysLarge toys)             	& 0.0699          & 0.139                	\\
%\hline
%\end{tabular}
%\caption{
%Summary of the expected rate-only and rate+shape p-values for 
%the Run 1-6 statistic, using \NumToys (up) and \NumToysLarge (bottom) toys for fake datasets (both $\beta=0$ and $\beta=1$).
%The Asimov p-value is calculated using the Asimov dataset while the mean one
%from \NumToys (up) and \NumToysLarge (bottom) toys and taking the mean 
%of the p-value distribution.
%}
%\label{tab:summary_sensi_manytoys_4p0108E20POT}
%\end{table}


%\begin{table}
%\centering
%\begin{tabular}{| c | c | c |}
%\hline
%P-value	                                          	& Asimov    	& Mean 		    	\\
%\hline
%Rate-only (\NumToys toys)        			& 0.0343           & 0.1004                   	\\
%Rate+shape (\NumToys toys)             		& 0.0249          	& 0.0649                 	\\
%\hline
%Rate-only (\NumToysLarge toys)        			& 0.0331           & 0.1007                   	\\
%Rate+shape (\NumToysLarge toys)             		& 0.0243          	& 0.0660                 	\\
%\hline
%\end{tabular}
%\caption{
%Summary of the expected rate-only and rate+shape p-values for 
%$8 \times 10^{20}$ POT, using \NumToys (up) and \NumToysLarge (bottom) toys for fake datasets (both $\beta=0$ and $\beta=1$).
%The Asimov p-value is calculated using the Asimov dataset while the mean one
%from \NumToys (up) and \NumToysLarge (bottom) toys and taking the mean 
%of the p-value distribution.
%}
%\label{tab:summary_sensi_manytoys_8E20POT}
%\end{table}

%





\begin{table}
\centering
\begin{tabular}{| c | c | c |}
\hline
		                                          	& p-value $\leq 0.10$    		& p-value $\leq 0.05$ 		\\
\hline
Fraction of toys	(rate-only)			& XXX\%					& XXX\%				\\
Fraction of toys	(rate+shape) 		& 58.0\%					& 46.5\%				\\
\hline
\end{tabular}
\caption{
The probability to have a p-value smaller than 0.05 and 0.10 is shown.
\NumToysLarge toys for
the Run 1-6 statistic
are produced.
Then for each toy the p-value is computed as described in section~\ref{subsec:pvalue_mean}.
}
\label{tab:pvalue_mean_0p05_lowpot}
\end{table}

\begin{table}
\centering
\begin{tabular}{| c | c | c |}
\hline
\# obs events $>$ \# exp bkg		& p-value    	 & Fraction of toys    	\\
\hline
Rate-only                          			&   XXX	    	& XXX\%             	\\
Rate+shape                          		&   0.102   	& 86.9\%                 	\\
\hline
\end{tabular}
\caption{
Summary of the expected rate-only and rate+shape mean p-values for 
the Run 1-6 statistic.
%\FakePotLow POT.
The Asimov p-value is calculated using the Asimov dataset while the 
mean p-value is calculated making \NumToysLarge toys and taking the mean 
of the p-value distribution.
Only fake datasets with number of observed events higher than the
expected background are considered.
The fraction of selected toys is shown as well. 
}
\label{tab:pvalue_mean_bkg_lowpot}
\end{table}



\begin{table}
\centering
\begin{tabular}{| c | c |}
\hline
p-value (rate-only) $>$ p-value (rate+shape)		& Probability          	\\
\hline
All toys									& XXX\%		 							\\
Toys w/ \# obs. events $>$ \# exp. bkg 			& XXX\%									\\
\hline
\end{tabular}
\caption{
The probability to have a rate-only p-value larger than the rate+shape p-value is shown.
\NumToysLarge toys for 
the Run 1-6 statistic.
%\FakePotLow POT are produced.
Then for each toy the p-value is computed as described in section~\ref{subsec:pvalue_mean}.
Both all toys are considered as well as the toys with the number of observed events larger
than the number of expected background.
}
\label{tab:pvalue_mean_rate_large_shape_lowpot}
\end{table}


